\documentclass[aspectratio=169,11pt]{beamer}

% ============================================================
% Presentación de defensa de tesis
% Efecto de la dilución aleatoria en el aprendizaje automático
% profundo con redes neuronales feed-forward
% Autor: Gamalerio Joaquín Ignacio
% ============================================================

\usetheme{Madrid}
\usecolortheme{default}
\setbeamertemplate{navigation symbols}{}
\setbeamertemplate{footline}[frame number]
\usepackage[spanish,es-nodecimaldot]{babel}
\usepackage[utf8]{inputenc}
\usepackage[T1]{fontenc}
\usepackage{lmodern}
\usepackage{amsmath,amssymb}
\usepackage{booktabs}
\usepackage{array}
\usepackage{tabularx}
\usepackage{graphicx}
\usepackage{grffile}
\usepackage{tikz}
\usepackage{xcolor}
\usepackage{ragged2e}
\usepackage{multicol}
\usepackage{hyperref}

\definecolor{famafblue}{HTML}{245B91}
\definecolor{boostred}{HTML}{C7565C}
\definecolor{softgray}{HTML}{F5F6F7}
\definecolor{darkgray}{HTML}{404040}

\setbeamercolor{normal text}{fg=darkgray,bg=white}
\setbeamercolor{alerted text}{fg=boostred}
\setbeamercolor{frametitle}{fg=white,bg=famafblue}
\setbeamercolor{progress bar}{fg=boostred,bg=softgray}
\setbeamercolor{title separator}{fg=boostred}
\setbeamertemplate{itemize item}{\raise0.2ex\hbox{\scriptsize$\blacktriangleright$}}
\setbeamertemplate{itemize subitem}{\raise0.2ex\hbox{\scriptsize$\circ$}}
\setbeamertemplate{caption}{\raggedright\insertcaption\par}

% Ruta base opcional. Si en Overleaf subís las imágenes respetando los paths
% originales de la tesis, no hace falta cambiar nada.
\graphicspath{{./}}

% Inserta imagen si existe; si no, deja un placeholder con el path exacto.
% Esto permite compilar la presentación aun antes de subir las imágenes.
\newcommand{\thesisimage}[3][]{%
  \IfFileExists{#2}{\includegraphics[#1]{#2}}{%
    \begin{center}
    \fbox{\begin{minipage}[c][#3][c]{0.92\linewidth}
      \centering\small\textbf{Imagen pendiente}
    \end{minipage}}
    \end{center}}
}

\newcommand{\keymsg}[1]{\begin{block}{Mensaje central}\large #1\end{block}}
\newcommand{\smallsource}[1]{}
\newcommand{\model}[1]{\textit{#1}}
\newcommand{\bd}{\model{BoostDropout}}
\newcommand{\bdn}{\model{BoostDropoutNet}}
\newcommand{\ofn}{\model{OverFitNet}}
\newcommand{\don}{\model{DropoutNet}}
\newcommand{\dcn}{\model{DropConnectNet}}

\title{Efecto de la dilución aleatoria en el aprendizaje automático profundo con redes neuronales \textit{feed-forward}}
\subtitle{Regularización estocástica mediante Dropout, DropConnect y BoostDropout}
\author{Gamalerio Joaquín Ignacio}
\institute{Facultad de Matemática, Astronomía, Física y Computación\\Universidad Nacional de Córdoba\\Director: Dr. Francisco A. Tamarit}
\date{Mayo, 2026}

\begin{document}

% ------------------------------------------------------------
\begin{frame}[plain]
  \titlepage
\end{frame}

% ------------------------------------------------------------

% ------------------------------------------------------------
%\begin{frame}{Mapa narrativo de la defensa}
%\begin{enumerate}
%  \item \textbf{Problema:} las redes con alta capacidad expresiva pueden memorizar y no generalizar.
%  \item \textbf{Herramientas clásicas:} Dropout y DropConnect introducen dilución estocástica binaria.
%  \item \textbf{Propuesta:} \bd{} reemplaza la supresión binaria por modulación multiplicativa de activaciones.
%  \item \textbf{Evaluación:} comparación controlada en MNIST con varias semillas.
%  \item \textbf{Resultado:} \bdn{} mejora la generalización en el protocolo estudiado.
%\end{enumerate}
%\vfill
%\begin{alertblock}{Alcance}
%No se busca demostrar superioridad universal, sino evaluar una idea nueva en un escenario controlado %\end{alertblock}
%\end{frame}

% ------------------------------------------------------------
\section{Motivación y marco conceptual}

\begin{frame}{Motivación: Queremos que la red aprenda, no que memorice.}
\begin{itemize}
  \item Una red neuronal aprende mirando siempre ejemplos conocidos.
  \item El objetivo no es solo que acierte esos ejemplos, sino que también funcione bien con ejemplos nuevos, que no ve durante el aprendizaje.
  \item A veces la red ``memoriza'' detalles del entrenamiento en lugar de aprender patrones útiles.
  \item Cuando eso ocurre, su desempeño baja frente a datos que nunca vio.
  \item Este problema se llama \textbf{sobreajuste}.
  \item La tesis estudia cómo reducirlo induciendo perturbaciones en la red durante el aprendizaje.
\end{itemize}

\vspace{0.25cm}
\end{frame}

% ------------------------------------------------------------
\begin{frame}{Objetivos de la tesis}
\begin{itemize}
  \item Estudiar cómo mejorar el desempeño de redes neuronales frente a datos nuevos.
  \item Comparar una red base con dos métodos conocidos: \textit{Dropout} y \textit{DropConnect}.
  \item Proponer y evaluar \bd{}, una variante que modifica la intensidad de algunas señales internas.
  \item Analizar si estas perturbaciones ayudan a reducir el \textbf{sobreajuste}.
  \item Usar un autoencoder para observar las representaciones aprendidas y estudiar su \textbf{esparsidad}.
\end{itemize}
\end{frame}

% ------------------------------------------------------------
% ------------------------------------------------------------
\begin{frame}{De una neurona a una red \textit{feed-forward}}
\begin{columns}[T,totalwidth=\textwidth]
\column{0.48\textwidth}
\centering
\textbf{Perceptrón simple}\par\vspace{0.1cm}
\thesisimage[width=\linewidth,height=5.4cm,keepaspectratio]{01_perceptron_simple.jpg}{4.8cm}

\column{0.48\textwidth}
\centering
\textbf{Red \textit{feed-forward}}\par\vspace{0.1cm}
\thesisimage[width=\linewidth,height=5.4cm,keepaspectratio]{02_perceptron_multicapa.jpg}{4.8cm}
\end{columns}

\vspace{0.25cm}
\begin{block}{Idea}
Una neurona transforma información de entrada en una señal de salida.
\end{block}
\note{
\begin{itemize}
  \item Una neurona artificial recibe varios números de entrada.
  \item A cada entrada le asigna una importancia, llamada peso.
  \item Combina esas entradas y produce una salida.
  \item Esa salida puede interpretarse como una decisión o predicción del modelo.
  \item Una red combina muchas neuronas organizadas en capas.
  \item La información avanza en una dirección: entrada, capas intermedias y salida.
  \item Cada neurona combina sus entradas y luego aplica una función de activación.
  \item La predicción final se obtiene a partir de la última capa.
\end{itemize}
}
\end{frame}

% ------------------------------------------------------------
\begin{frame}{Diagrama de entrenamiento}
\centering
\resizebox{0.82\textwidth}{!}{%
\begin{tikzpicture}[node distance=1.45cm, every node/.style={font=\small}]
\node[draw, rounded corners, fill=softgray, minimum width=2.3cm, minimum height=0.85cm] (train) {Entrenamiento};
\node[draw, rounded corners, fill=softgray, right of=train, xshift=1.55cm, minimum width=2.2cm, minimum height=0.85cm] (pred) {Predicción};
\node[draw, rounded corners, fill=softgray, right of=pred, xshift=1.55cm, minimum width=2.0cm, minimum height=0.85cm] (loss) {Pérdida};
\node[draw, rounded corners, fill=softgray, below of=loss, yshift=-0.15cm, minimum width=2.4cm, minimum height=0.85cm] (grad) {Gradientes};
\node[draw, rounded corners, fill=softgray, left of=grad, xshift=-1.55cm, minimum width=2.5cm, minimum height=0.85cm] (upd) {Actualizar pesos};

\node[draw, rounded corners, fill=softgray, below of=train, yshift=-1.20cm, minimum width=2.3cm, minimum height=0.85cm] (val) {Validación};
\node[draw, rounded corners, fill=softgray, right of=val, xshift=1.55cm, minimum width=2.2cm, minimum height=0.85cm] (sel) {Elegir modelo};
\node[draw, rounded corners, fill=softgray, right of=sel, xshift=1.55cm, minimum width=2.0cm, minimum height=0.85cm] (test) {Test final};

\draw[->, thick] (train) -- (pred);
\draw[->, thick] (pred) -- (loss);
\draw[->, thick] (loss) -- (grad);
\draw[->, thick] (grad) -- (upd);
\draw[->, thick] (upd.west) to[out=180,in=-90] (train.south);

\draw[->, thick] (val) -- (sel);
\draw[->, thick] (sel) -- (test);
\draw[->, thick, dashed] (train) -- (val);
\end{tikzpicture}
}

\vspace{0.2cm}
\begin{block}{Resumen}
Entrenamiento ajusta los pesos. Validación ayuda a elegir hiperparámetros de entrenamiento.
Test se usa solo al final para evaluar datos no vistos.
\end{block}
\note{
\begin{itemize}
  \item La red aprende usando un conjunto de entrenamiento.
  \item Para cada ejemplo, hace una predicción y la compara con la respuesta correcta.
  \item Esa comparación se mide con una función de pérdida.
  \item La retropropagación calcula gradientes.
  \item El optimizador, por ejemplo descenso por el gradiente o Adam, actualiza los pesos para reducir la pérdida.
  \item El conjunto de validación se usa durante el entrenamiento para elegir hiperparámetros y seleccionar modelos.
  \item El conjunto de test se reserva hasta el final para medir el desempeño definitivo.
\end{itemize}
}
\end{frame}

% ------------------------------------------------------------
%\begin{frame}{Clasificación supervisada en MNIST}
%\begin{columns}[T,totalwidth=\textwidth]
%\column{0.42\textwidth}
%\begin{itemize}
%  \item Como ejemplo de tarea, usamos MNIST: imágenes de dígitos manuscritos.
%  \item La entrada es una imagen de $28\times28$ píxeles.
 % \item La salida esperada es una clase entre $0$ y $9$.
  %\item La red asigna un puntaje a cada dígito posible.
  %\item La predicción final es el dígito con mayor puntaje.
%\end{itemize}

%\column{0.45\textwidth}
%\thesisimage[width=\linewidth]{mnist_ejemplos.jpg}{4.4cm}
%\end{columns}
%\end{frame}

% ------------------------------------------------------------
\begin{frame}{Sobreajuste y generalización}
\small
\begin{columns}[T,totalwidth=\textwidth]
\column{0.35\textwidth}
\vspace{0.25cm}
Sobreajuste:
\begin{itemize}
  \item Hay sobreajuste cuando el \textbf{error de test} es muy alto comparado con el \textbf{error de entrenamiento}
\end{itemize}
Diagnóstico temprano:
\begin{itemize}
  \item La \textbf{pérdida de entrenamiento} puede seguir bajando.
  \item La \textbf{pérdida de validación} puede dejar de mejorar y comenzar a subir.
\end{itemize}
\column{0.62\textwidth}
\centering
\thesisimage[width=\linewidth,height=0.72\textheight,keepaspectratio]{15_curvas_loss_overfitting.jpg}{5.4cm}
\end{columns}
\smallsource{15_curvas_loss_overfitting.jpg}

\note{
\begin{itemize}
  \item Explicar qué significa error de entrenamiento y test, ya que usaremos esa métrica en los resultados
  \item Remarcar que el sobreajuste aparece cuando el error de entrenamiento es bajo, pero el error de test es alto.
  \item Usar la figura para explicar la pérdida de entrenamiento: puede seguir bajando porque la red ajusta cada vez mejor los datos conocidos.
  \item Usar la figura para explicar la pérdida de validación: si deja de bajar y empieza a subir, indica que el modelo empieza a perder capacidad de generalizar.
\end{itemize}
}
\end{frame}

% ------------------------------------------------------------
\begin{frame}{Regularización}
\begin{center}
{\fontsize{28}{32}\bfseries ¿Qué es la regularización?}

\vspace{0.65cm}
{\fontsize{18}{22}\bfseries ¿Qué entendemos por dilución y perturbación en el entrenamiento?}
\end{center}
\note{
\begin{itemize}
  \item Regularizar es modificar el aprendizaje para reducir el error de test, aunque eso no implique reducir el error de entrenamiento.
  \item Usar como disparadores las ideas de dilución, penalización y perturbación.
  \item Sin regularización, la red tiene mayor libertad efectiva y más riesgo de depender de patrones frágiles.
  \item Ese exceso de libertad puede aumentar el sobreajuste.
  \item Con regularización estocástica, perturbamos la propagación durante entrenamiento.
  \item La idea es obligar a la red a distribuir información de manera más robusta.
  \item El objetivo final es mejorar el desempeño en datos nuevos, no necesariamente minimizar más el error de entrenamiento.
  \item La dilución es una perturbación binaria.
\end{itemize}
}
\end{frame}

% ------------------------------------------------------------
\section{Métodos de regularización}

\begin{frame}{Dropout y DropConnect: dilución aleatoria}
\begin{columns}[T,totalwidth=\textwidth]
\column{0.48\textwidth}
\centering
\textbf{Dropout}\par\vspace{0.1cm}
\thesisimage[width=\linewidth,height=0.58\textheight,keepaspectratio,trim=0 0 0 1.25cm,clip]{dropconnect y dropout/10_inverted_dropout_training.jpg}{4.4cm}

\column{0.48\textwidth}
\centering
\textbf{DropConnect}\par\vspace{0.1cm}
\thesisimage[width=\linewidth,height=0.58\textheight,keepaspectratio,trim=0 0 0 1.25cm,clip]{dropconnect y dropout/12_dropconnect_training.jpg}{4.4cm}
\end{columns}

\vspace{0.25cm}
\begin{block}{Idea central}
Dropout y DropConnect son dos estrategias de dilución aleatoria durante el entrenamiento.
\end{block}
\note{
\begin{itemize}
  \item En Dropout, en cada iteración del entrenamiento, una fracción de activaciones se anula aleatoriamente.
  \item La perturbación de Dropout es binaria: cada activación queda activa o apagada.
  \item Con Dropout, la red no puede depender rígidamente de unidades específicas.
  \item Puede interpretarse como entrenamiento de muchas subredes implícitas con pesos compartidos.
  \item En DropConnect, la máscara aleatoria actúa sobre pesos individuales, no sobre neuronas completas.
  \item En DropConnect, la activación puede existir, pero algunas conexiones salientes o entrantes se anulan.
  \item La conectividad efectiva cambia en cada iteración.
  \item DropConnect también es una perturbación binaria.
  \item Ambos métodos implementan normalizacion durante el entrenamiento
\end{itemize}
}
\end{frame}

% ------------------------------------------------------------
\begin{frame}{\bd: idea central}
\begin{columns}[T,totalwidth=\textwidth]
\column{0.40\textwidth}
\vspace{0.45cm}
\begin{block}{Máscara propuesta}
\small
\[
 m_j = \begin{cases}
 1+\lambda, & \text{con probabilidad } p,\\
 1, & \text{con probabilidad } 1-p,
 \end{cases}
\]
\vspace{-0.15cm}
\[
 \widetilde h_j = m_j h_j
\]
\end{block}
\column{0.64\textwidth}
\vfill
\begin{center}
\thesisimage[width=\linewidth,height=0.76\textheight,keepaspectratio]{14_boostdropout_training.jpg}{5.8cm}
\end{center}
\smallsource{14_boostdropout_training.jpg}
\end{columns}
\note{
\begin{itemize}
  \item No apaga necesariamente la señal.
  \item Modula multiplicativamente la activación.
  \item Reemplaza ``apagar/conservar'' por ``modificar amplitud''.
  \item Introduce dos hiperparámetros principales: frecuencia $p$ e intensidad $\lambda$.
  \item El valor esperado de la máscara es $\mathbb{E}[m_j] = 1 + p\lambda$.
  \item Si se quiere preservar la escala esperada, puede usarse $\widehat m_j = m_j/(1+p\lambda)$, de modo que $\mathbb{E}[\widehat m_j]=1$.
  \item Si $\lambda=-1$, la máscara pasa a ser binaria sobre $\{0,1\}$.
  \item Con normalización adecuada, ese caso recupera la lógica de \textit{inverted dropout}.
  \item En esta tesis, la configuración seleccionada de \bd{} usa $p=0.7$, $\lambda=0.9$ y máscara no normalizada.
  \item Antes de avanzar al diseño experimental, conviene recapitular los tres métodos como una misma familia de regularización estocástica.
  \item Dropout actúa sobre activaciones con una perturbación binaria: suprime aleatoriamente unidades.
  \item DropConnect actúa sobre conexiones con una perturbación binaria: suprime aleatoriamente pesos.
  \item \bd{} actúa sobre activaciones con una perturbación no binaria y multiplicativa: modula estocásticamente la amplitud.
  \item La diferencia central está en el nivel donde actúa cada máscara y en si la perturbación es binaria o multiplicativa.
\end{itemize}
}
\end{frame}

% ------------------------------------------------------------
\section{Diseño experimental}

\begin{frame}{Experimentación Principal}
\centering
\vspace{0.15cm}
{\Large\bfseries Comparación de métodos de regularización en redes neuronales \textit{feed-forward}}

\vspace{0.35cm}
\begin{columns}[T,totalwidth=\textwidth]
\column{0.47\textwidth}
\thesisimage[width=\linewidth,height=0.59\textheight,keepaspectratio]{mnist_ejemplos.jpg}{4.8cm}
\smallsource{mnist_ejemplos.jpg}

\column{0.50\textwidth}
\begin{itemize}
  \item MNIST como tarea de clasificación supervisada.
  \item Arquitectura compartida para todas las redes comparadas.
  \item Fase 1: exploración de hiperparámetros y selección por validación.
  \item Fase 2: evaluación final multisemilla con configuraciones fijadas.
\end{itemize}
\end{columns}

\vspace{0.35cm}
\centering
%\begin{tabular}{|l|c|c|c|c|}
%\hline
%\textbf{Modelo} & \ofn & \don & \dcn & \bdn \\
%\hline
%\textbf{Regularización} & Sin reg. & Dropout & DropConnect & \bd \\
%\hline
%\end{tabular}
\note{
\begin{itemize}
  \item Entrenar redes \textit{feed-forward} sobre MNIST.
  \item Comparar error de test entre \ofn{}, \don{}, \dcn{} y \bdn{}.
  \item Evaluar si \bd{} mejora la generalización frente a regularizadores clásicos.
  \item La diapositiva introduce solo la experimentación principal: clasificación con redes densas \textit{feed-forward}.
  \item No introducir todavía el análisis cualitativo con autoencoders.
  \item Fase 1: exploración mediante grid search de hiperparámetros, varias semillas y selección por validación.
  \item Fase 2: evaluación final con configuraciones fijadas, entrenamiento multisemilla y reporte en validación y test.
  \item Remarcar que la evaluación multisemilla evita que una corrida afortunada defina la conclusión.
  \item Los modelos comparados son \ofn{}, \don{}, \dcn{} y \bdn{}.
  \item \ofn{} es la arquitectura base sin regularización estocástica.
  \item \don{} usa la misma arquitectura con Dropout en capas ocultas.
  \item \dcn{} usa la misma arquitectura con DropConnect en capas densas.
  \item \bdn{} usa la misma arquitectura con modulación multiplicativa \bd{}.
  \item La comparación se concentra en el mecanismo de regularización, no en cambiar tamaño, profundidad o dataset.
  \item MNIST contiene dígitos manuscritos de $0$ a $9$, representados como imágenes de $28\times28$ píxeles.
  \item Se usa un subconjunto estratificado para reducir costo computacional y volver visible el sobreajuste.
  \item La decisión metodológica no busca maximizar rendimiento absoluto en MNIST, sino estudiar regularización en un escenario controlado.
  \item La arquitectura base es una red densa \textit{feed-forward}, con capas ocultas ReLU y salida de 10 clases.
  \item La misma arquitectura se mantiene para los cuatro modelos comparados.
\end{itemize}
}
\end{frame}

\begin{frame}{Criterios de selección: MMV y \textit{early stopping}}
\begin{columns}[T,totalwidth=\textwidth]
\column{0.48\textwidth}
\textbf{MMV}
\begin{itemize}
  \item Mejor modelo según pérdida de validación.
  \item Selecciona el modelo que alcanza la menor pérdida de validación durante todo el entrenamiento.
\end{itemize}

\column{0.48\textwidth}
\textbf{\textit{Early stopping}}
\begin{itemize}
  \item Modelo elegido por una política de detención temprana.
  \item Usa pérdida de validación como métrica y una tolerancia definida.
\end{itemize}
\end{columns}
\note{
Comparar ambos criterios permite verificar si la conclusión depende del modo de seleccionar el modelo final.
}
\end{frame}

% ------------------------------------------------------------
\section{Resultados principales}

\begin{frame}{Configuraciones seleccionadas}
\textbf{Hiperparámetros compartidos}
\vspace{-0.1cm}
\begin{table}
\centering
\scriptsize
\begin{tabularx}{0.95\textwidth}{lX}
\toprule
Decisión & Uso en el experimento \\
\midrule
Optimizador & Adam. \\
Batch size & $256$. \\
Activación & ReLU en capas ocultas. \\
Pérdida & Entropía cruzada para clasificación multiclase. \\
Validación & Selección de hiperparámetros y artefactos. \\
Test & Reservado para evaluación final. \\
\bottomrule
\end{tabularx}
\end{table}

\vspace{-0.15cm}
\textbf{Hiperparámetros de la búsqueda en grilla}
\vspace{-0.1cm}
\begin{table}
\centering
\scriptsize
\begin{tabularx}{0.98\textwidth}{lXXX}
\toprule
Modelo & $p$ explorado & $\lambda$ explorado & Normalización explorada \\
\midrule
\don & 0.1, 0.3, \alert{0.5}, 0.7, 0.9 & -- & -- \\
\dcn & 0.1, 0.3, \alert{0.5}, 0.7, 0.9 & -- & -- \\
\bdn & 0.1, 0.3, 0.5, \alert{0.7}, 0.9 & -0.6, -0.3, 0.0, 0.3, 0.6, \alert{0.9} & True, \alert{False} \\
\bottomrule
\end{tabularx}
\end{table}
\note{
\begin{itemize}
  \item Separación estricta: el test no interviene en la búsqueda de hiperparámetros ni en la selección de modelos.
  \item Mensaje central: la configuración seleccionada de \bd{} corresponde a una modulación frecuente, positiva y no normalizada.
\end{itemize}
}
\end{frame}

% ------------------------------------------------------------
\begin{frame}{Resultados principales}
\begin{center}
\vspace{1.6cm}
{\fontsize{28}{34}\bfseries Resultados de experimentación Principal}
\end{center}
\end{frame}

% ------------------------------------------------------------
\begin{frame}{Resultado global: criterio MMV}
\vspace{0.15cm}
\begin{columns}[c,totalwidth=\textwidth]
\column{0.50\textwidth}
\centering
\textbf{Error de test}\par\vspace{0.05cm}
\thesisimage[width=\linewidth,height=0.72\textheight,keepaspectratio]{imagenes_resultados/entrenamiento_principal/best_model/boxplot_error_test_best_model.jpg}{6.0cm}
\column{0.50\textwidth}
\centering
\textbf{Pérdida de test}\par\vspace{0.05cm}
\thesisimage[width=\linewidth,height=0.72\textheight,keepaspectratio]{imagenes_resultados/entrenamiento_principal/best_model/boxplot_perdida_test_best_model.jpg}{6.0cm}
\end{columns}
% \begin{table}
% \centering
% \small
% \begin{tabular}{lcc}
% \toprule
% Modelo & Error test & Pérdida test \\
% \midrule
% \ofn & 0.0687 & 0.2311 \\
% \don & 0.0609 & 0.2129 \\
% \dcn & 0.0614 & 0.2133 \\
% \bdn & \alert{0.0480} & \alert{0.1598} \\
% \bottomrule
% \end{tabular}
% \end{table}
% \begin{block}{Lectura}
% Bajo MMV, \bdn{} obtiene menor error y menor pérdida media de test.
% \end{block}
\smallsource{imagenes_resultados/entrenamiento_principal/best_model/boxplot_error_test_best_model.jpg}
\note{
\begin{itemize}
  \item \ofn{} presenta el peor desempeño: el sobreajuste afecta la generalización.
  \item \don{} y \dcn{} mejoran respecto de la red base.
  \item Ambos regularizadores clásicos quedan muy cercanos entre sí.
  \item \bdn{} desplaza la distribución hacia valores menores de error y pérdida.
  \item La mejora de \bdn{} no aparece solo en un valor medio aislado: también se observa en medianas, boxplots y métricas de validación.
\end{itemize}
}
\end{frame}

% ------------------------------------------------------------
\begin{frame}{Resultado global: \textit{early stopping}}
\vspace{0.15cm}
\begin{columns}[c,totalwidth=\textwidth]
\column{0.50\textwidth}
\centering
\textbf{Error de test}\par\vspace{0.05cm}
\thesisimage[width=\linewidth,height=0.78\textheight,keepaspectratio]{imagenes_resultados/entrenamiento_principal/early_stopping/boxplot_error_test_early_stopping.jpg}{6.0cm}
\column{0.50\textwidth}
\centering
\textbf{Pérdida de test}\par\vspace{0.05cm}
\thesisimage[width=\linewidth,height=0.78\textheight,keepaspectratio]{imagenes_resultados/entrenamiento_principal/early_stopping/boxplot_perdida_test_early_stopping.jpg}{6.0cm}
\end{columns}
% \begin{table}
% \centering
% \small
% \begin{tabular}{lcc}
% \toprule
% Modelo & Error test & Pérdida test \\
% \midrule
% \ofn & 0.0687 & 0.2310 \\
% \don & 0.0609 & 0.2129 \\
% \dcn & 0.0614 & 0.2133 \\
% \bdn & \alert{0.0491} & \alert{0.1640} \\
% \bottomrule
% \end{tabular}
% \end{table}
\smallsource{imagenes_resultados/entrenamiento_principal/early_stopping/boxplot_error_test_early_stopping.jpg}
\note{
\begin{itemize}
  \item Comparar MMV y \textit{early stopping} permite verificar si la conclusión principal depende del criterio de selección.
  \item Errores de test con MMV: \ofn{} 0.0687, \don{} 0.0609, \dcn{} 0.0614, \bdn{} 0.0480.
  \item Errores de test con \textit{early stopping}: \ofn{} 0.0687, \don{} 0.0609, \dcn{} 0.0614, \bdn{} 0.0491.
  \item La conclusión principal no depende fuertemente del criterio de selección.
  \item Para \bdn{}, MMV obtiene un resultado levemente mejor.
  \item Esto sugiere que parte de su ventaja aparece en un régimen más tardío del entrenamiento.
\end{itemize}
}
\end{frame}

% ------------------------------------------------------------
\begin{frame}{Distribución de la mejor época}
\vspace{0.15cm}
\begin{columns}[c,totalwidth=\textwidth]
\column{0.50\textwidth}
\centering
\textbf{Mejor época MMV}\par\vspace{0.05cm}
\thesisimage[width=\linewidth,height=0.78\textheight,keepaspectratio]{imagenes_resultados/entrenamiento_principal/best_model/boxplot_mejor_epoca_best_model.jpg}{6.0cm}
\column{0.50\textwidth}
\centering
\textbf{Mejor época \textit{early stopping}}\par\vspace{0.05cm}
\thesisimage[width=\linewidth,height=0.78\textheight,keepaspectratio]{imagenes_resultados/entrenamiento_principal/early_stopping/boxplot_mejor_epoca_early_stopping.jpg}{6.0cm}
\end{columns}
% \begin{table}
% \centering
% \small
% \begin{tabular}{lcc}
% \toprule
% Modelo & MMV & Early \\
% \midrule
% \ofn & 7.8 & 7.8 \\
% \don & 12.0 & 12.0 \\
% \dcn & 16.3 & 16.3 \\
% \bdn & \alert{43.8} & \alert{29.2} \\
% \bottomrule
% \end{tabular}
% \end{table}
\end{frame}

% ------------------------------------------------------------
\begin{frame}{Curvas representativas de validación (MMV)}
\vspace{0.1cm}
\begin{columns}[c,totalwidth=\textwidth]
\column{0.50\textwidth}
\centering
\textbf{Error de validación}\par\vspace{0.05cm}
\thesisimage[width=\linewidth,height=0.72\textheight,keepaspectratio]{imagenes_resultados/entrenamiento_principal/best_model/curva_error_validacion_best_model.jpg}{5.6cm}
\column{0.50\textwidth}
\centering
\textbf{Pérdida de validación}\par\vspace{0.05cm}
\thesisimage[width=\linewidth,height=0.72\textheight,keepaspectratio]{imagenes_resultados/entrenamiento_principal/best_model/curva_perdida_validacion_best_model.jpg}{5.6cm}
\end{columns}
\smallsource{imagenes_resultados/entrenamiento_principal/best_model/curva_perdida_validacion_best_model.jpg}
\note{
\begin{itemize}
  \item \ofn{} muestra deterioro tardío de validación.
  \item Dropout y DropConnect atenúan ese patrón.
  \item \bdn{} mantiene una validación más estable en épocas tardías.
\end{itemize}
}
\end{frame}

% ------------------------------------------------------------
\begin{frame}{Dependencia con $p$ (MMV)}
\vspace{0.1cm}
\begin{columns}[c,totalwidth=\textwidth]
\column{0.50\textwidth}
\centering
\textbf{Error de test}\par\vspace{0.05cm}
\thesisimage[width=\linewidth,height=0.72\textheight,keepaspectratio]{imagenes_resultados/graficos_variando_p/comparison_error_test_best_model.jpg}{5.6cm}
\column{0.50\textwidth}
\centering
\textbf{Pérdida de test}\par\vspace{0.05cm}
\thesisimage[width=\linewidth,height=0.72\textheight,keepaspectratio]{imagenes_resultados/graficos_variando_p/comparison_perdida_test_best_model.jpg}{5.6cm}
\end{columns}
\note{
\begin{itemize}
  \item $p$ controla la frecuencia de modulación o enmascaramiento.
  \item \bdn{} muestra una región favorable alrededor de $p=0.7$.
  \item La mejora no depende de un punto aislado: se mantiene en una zona amplia del rango explorado.
\end{itemize}
}
\end{frame}

% ------------------------------------------------------------
\begin{frame}{Dependencia con $p$ (\textit{early stopping})}
\vspace{0.1cm}
\begin{columns}[c,totalwidth=\textwidth]
\column{0.50\textwidth}
\centering
\textbf{Error de test}\par\vspace{0.05cm}
\thesisimage[width=\linewidth,height=0.72\textheight,keepaspectratio]{imagenes_resultados/graficos_variando_p/comparison_error_test_early_stopping.jpg}{5.6cm}
\column{0.50\textwidth}
\centering
\textbf{Pérdida de test}\par\vspace{0.05cm}
\thesisimage[width=\linewidth,height=0.72\textheight,keepaspectratio]{imagenes_resultados/graficos_variando_p/comparison_perdida_test_early_stopping.jpg}{5.6cm}
\end{columns}
\note{
\begin{itemize}
  \item Se repite el patrón observado con MMV.
  \item La región alrededor de $p=0.7$ sigue siendo favorable para \bdn{}.
\end{itemize}
}
\end{frame}


% ------------------------------------------------------------
\begin{frame}{Dependencia con $\lambda$ (MMV)}
\vspace{0.1cm}
\begin{columns}[c,totalwidth=\textwidth]
\column{0.50\textwidth}
\centering
\textbf{Error de test}\par\vspace{0.05cm}
\thesisimage[width=\linewidth,height=0.72\textheight,keepaspectratio]{imagenes_resultados/graficos variando lambda/comparison_lambda_error_test_best_model.jpg}{5.6cm}
\column{0.50\textwidth}
\centering
\textbf{Pérdida de test}\par\vspace{0.05cm}
\thesisimage[width=\linewidth,height=0.72\textheight,keepaspectratio]{imagenes_resultados/graficos variando lambda/comparison_lambda_perdida_test_best_model.jpg}{5.6cm}
\end{columns}
\note{
\begin{itemize}
  \item $\lambda$ controla la intensidad de la modulación multiplicativa.
  \item La región positiva contiene los mejores resultados de \bdn{}.
  \item Valores negativos, especialmente cerca de $\lambda=-1$ sin normalización, degradan fuertemente el desempeño.
\end{itemize}
}
\end{frame}

% ------------------------------------------------------------
\begin{frame}{Dependencia con $\lambda$ (\textit{early stopping})}
\vspace{0.1cm}
\begin{columns}[c,totalwidth=\textwidth]
\column{0.50\textwidth}
\centering
\textbf{Error de test}\par\vspace{0.05cm}
\thesisimage[width=\linewidth,height=0.72\textheight,keepaspectratio]{imagenes_resultados/graficos variando lambda/comparison_lambda_error_test_early_stopping.jpg}{5.6cm}
\column{0.50\textwidth}
\centering
\textbf{Pérdida de test}\par\vspace{0.05cm}
\thesisimage[width=\linewidth,height=0.72\textheight,keepaspectratio]{imagenes_resultados/graficos variando lambda/comparison_lambda_perdida_test_early_stopping.jpg}{5.6cm}
\end{columns}
\note{
\begin{itemize}
  \item Se repite el patrón observado con MMV.
  \item La región positiva sigue siendo la zona favorable para \bdn{}.
\end{itemize}
}
\end{frame}

% ------------------------------------------------------------
\begin{frame}{Conclusión del protocolo principal}
\large
\begin{enumerate}
  \item La regularización estocástica mejora sistemáticamente a la red base.
  \item Dropout y DropConnect reducen el sobreajuste y obtienen resultados similares entre sí.
  \item \bdn{} supera a las tres arquitecturas de referencia en error y pérdida de test.
  \item La ventaja se sostiene bajo MMV, \textit{early stopping}, curvas temporales y análisis de sensibilidad.
\end{enumerate}
\note{
La modulación multiplicativa no binaria aparece como una estrategia efectiva de regularización en el escenario experimental estudiado.
}
\end{frame}

% ------------------------------------------------------------
\section{Análisis cualitativo con autoencoders}


\begin{frame}{Experimentación secundaria: Autoencoders}
\begin{columns}[T,totalwidth=\textwidth]
\column{0.58\textwidth}
\centering
{\large\bfseries Esquema de Autoencoder}\par\vspace{0.08cm}
\thesisimage[width=\linewidth,height=0.82\textheight,keepaspectratio,trim=0 0 0 10mm,clip]{15_autoencoder_architecture.jpg}{6.3cm}
\smallsource{15_autoencoder_architecture.jpg}

\column{0.39\textwidth}
\begin{itemize}
  \item Entrenar autoencoders con MNIST.
  \item Comparar las mismas estrategias de regularización.
  \item Analizar esparsidad representacional y paramétrica.
  \item Observar características aprendidas.
\end{itemize}
\end{columns}
\note{
\begin{itemize}
  \item El protocolo principal mide desempeño predictivo: error y pérdida sobre datos de test.
  \item Para complementar esa evaluación, se inspeccionan las representaciones aprendidas por cada regularizador.
  \item Los autoencoders no reemplazan la evaluación de clasificación; sirven como herramienta exploratoria.
  \item Un autoencoder aprende a reconstruir su entrada.
  \item Permite observar representaciones internas y características aprendidas.
  \item La pregunta cualitativa es qué tipo de estructura interna induce cada mecanismo de regularización.
  \item El objetivo es mirar si los mecanismos de regularización producen firmas cualitativas distintas en las representaciones.
\end{itemize}
}
\end{frame}

% ------------------------------------------------------------


\begin{frame}{Resultados secundarios}
\begin{center}
\vspace{1.6cm}
{\fontsize{28}{34}\bfseries Resultados de experimentación secundaria con autoencoders}
\end{center}
\end{frame}

% ------------------------------------------------------------

% ------------------------------------------------------------
\begin{frame}{Visualización de coadaptaciones: autoencoder base y con Dropout}
\vspace{0.1cm}
\centering
\setlength{\tabcolsep}{0.03cm}
\renewcommand{\arraystretch}{0}
\begin{tabular}{cc}
\textbf{Coadaptaciones autoencoder base} & \textbf{Coadaptaciones Dropout} \\
\thesisimage[width=0.50\textwidth,height=0.78\textheight,keepaspectratio]{imagenes_resultados/autoencoder/OverFitNet/coadaptation.jpg}{6.0cm} &
\thesisimage[width=0.50\textwidth,height=0.78\textheight,keepaspectratio]{imagenes_resultados/autoencoder/Dropout/coadaptation.jpg}{6.0cm}
\end{tabular}
\smallsource{imagenes_resultados/autoencoder/OverFitNet/coadaptation.jpg; imagenes_resultados/autoencoder/Dropout/coadaptation.jpg}
\end{frame}

% ------------------------------------------------------------
\begin{frame}{Visualización de coadaptaciones: con DropConnect y BoostDropout}
\vspace{0.1cm}
\centering
\setlength{\tabcolsep}{0.03cm}
\renewcommand{\arraystretch}{0}
\begin{tabular}{cc}
\textbf{Coadaptaciones DropConnect} & \textbf{Coadaptaciones BoostDropout} \\
\thesisimage[width=0.50\textwidth,height=0.78\textheight,keepaspectratio]{imagenes_resultados/autoencoder/DropConnect/coadaptation.jpg}{6.0cm} &
\thesisimage[width=0.50\textwidth,height=0.78\textheight,keepaspectratio]{imagenes_resultados/autoencoder/BoostDropout/coadaptation.jpg}{6.0cm}
\end{tabular}
\smallsource{imagenes_resultados/autoencoder/DropConnect/coadaptation.jpg; imagenes_resultados/autoencoder/BoostDropout/coadaptation.jpg}
\end{frame}

% ------------------------------------------------------------
\begin{frame}{Decodificaciones latentes: autoencoder base y con Dropout}
\vspace{0.1cm}
\centering
\setlength{\tabcolsep}{0.03cm}
\renewcommand{\arraystretch}{0}
\begin{tabular}{cc}
\textbf{Decodificaciones autoencoder base} & \textbf{Decodificaciones Dropout} \\
\thesisimage[width=0.50\textwidth,height=0.78\textheight,keepaspectratio,trim=0 0 0 35bp,clip]{imagenes_resultados/autoencoder/OverFitNet/single_latent_decodes.jpg}{6.0cm} &
\thesisimage[width=0.50\textwidth,height=0.78\textheight,keepaspectratio,trim=0 0 0 35bp,clip]{imagenes_resultados/autoencoder/Dropout/single_latent_decodes.jpg}{6.0cm}
\end{tabular}
\smallsource{imagenes_resultados/autoencoder/OverFitNet/single_latent_decodes.jpg; imagenes_resultados/autoencoder/Dropout/single_latent_decodes.jpg}
\end{frame}

% ------------------------------------------------------------
\begin{frame}{Decodificaciones latentes: con DropConnect y BoostDropout}
\vspace{0.1cm}
\centering
\setlength{\tabcolsep}{0.03cm}
\renewcommand{\arraystretch}{0}
\begin{tabular}{cc}
\textbf{Decodificaciones DropConnect} & \textbf{Decodificaciones BoostDropout} \\
\thesisimage[width=0.50\textwidth,height=0.78\textheight,keepaspectratio,trim=0 0 0 35bp,clip]{imagenes_resultados/autoencoder/DropConnect/single_latent_decodes.jpg}{6.0cm} &
\thesisimage[width=0.50\textwidth,height=0.78\textheight,keepaspectratio,trim=0 0 0 35bp,clip]{imagenes_resultados/autoencoder/BoostDropout/single_latent_decodes.jpg}{6.0cm}
\end{tabular}
\smallsource{imagenes_resultados/autoencoder/DropConnect/single_latent_decodes.jpg; imagenes_resultados/autoencoder/BoostDropout/single_latent_decodes.jpg}
\note{
\begin{itemize}
  \item Dropout tiende a mostrar detectores más localizados y menor coadaptación visual.
  \item DropConnect aparece en una posición intermedia.
  \item \bd{} no replica exactamente la firma visual de Dropout.
  \item Esto sugiere que la mejora de \bdn{} no debe interpretarse como una versión más intensa de Dropout.
  \item \bd{} regulariza, pero probablemente mediante un mecanismo cualitativamente distinto al de la supresión binaria de activaciones.
\end{itemize}
}
\end{frame}

% ------------------------------------------------------------
\subsection{Esparsidad representacional}

\begin{frame}{Esparsidad representacional: autoencoder base y con Dropout}
\vspace{0.1cm}
\centering
\setlength{\tabcolsep}{0.03cm}
\renewcommand{\arraystretch}{0}
\begin{tabular}{cc}
\textbf{Esparsidad autoencoder base} & \textbf{Esparsidad Dropout} \\
\thesisimage[width=0.49\textwidth,height=0.78\textheight,keepaspectratio]{imagenes_resultados/autoencoder/OverFitNet/sparsity.jpg}{6.0cm} &
\thesisimage[width=0.49\textwidth,height=0.78\textheight,keepaspectratio]{imagenes_resultados/autoencoder/Dropout/sparsity.jpg}{6.0cm}
\end{tabular}
\smallsource{imagenes_resultados/autoencoder/OverFitNet/sparsity.jpg; imagenes_resultados/autoencoder/Dropout/sparsity.jpg}
\end{frame}

% ------------------------------------------------------------
\begin{frame}{Esparsidad representacional: con DropConnect y BoostDropout}
\vspace{0.1cm}
\centering
\setlength{\tabcolsep}{0.03cm}
\renewcommand{\arraystretch}{0}
\begin{tabular}{cc}
\textbf{Esparsidad DropConnect} & \textbf{Esparsidad BoostDropout} \\
\thesisimage[width=0.49\textwidth,height=0.78\textheight,keepaspectratio]{imagenes_resultados/autoencoder/DropConnect/sparsity.jpg}{6.0cm} &
\thesisimage[width=0.49\textwidth,height=0.78\textheight,keepaspectratio]{imagenes_resultados/autoencoder/BoostDropout/sparsity.jpg}{6.0cm}
\end{tabular}
\smallsource{imagenes_resultados/autoencoder/DropConnect/sparsity.jpg; imagenes_resultados/autoencoder/BoostDropout/sparsity.jpg}
\end{frame}

% ------------------------------------------------------------
\subsection{Esparsidad paramétrica}

\begin{frame}{Esparsidad paramétrica: autoencoder base y con Dropout}
\vspace{0.1cm}
\centering
\setlength{\tabcolsep}{0.03cm}
\renewcommand{\arraystretch}{0}
\begin{tabular}{cc}
\textbf{Esparsidad autoencoder base} & \textbf{Esparsidad Dropout} \\
\thesisimage[width=0.49\textwidth,height=0.78\textheight,keepaspectratio]{imagenes_resultados/autoencoder/OverFitNet/first_layer_parameter_sparsity.jpg}{6.0cm} &
\thesisimage[width=0.49\textwidth,height=0.78\textheight,keepaspectratio]{imagenes_resultados/autoencoder/Dropout/first_layer_parameter_sparsity.jpg}{6.0cm}
\end{tabular}
\smallsource{imagenes_resultados/autoencoder/OverFitNet/first_layer_parameter_sparsity.jpg; imagenes_resultados/autoencoder/Dropout/first_layer_parameter_sparsity.jpg}
\end{frame}

% ------------------------------------------------------------
\begin{frame}{Esparsidad paramétrica: con DropConnect y BoostDropout}
\vspace{0.1cm}
\centering
\setlength{\tabcolsep}{0.03cm}
\renewcommand{\arraystretch}{0}
\begin{tabular}{cc}
\textbf{Esparsidad DropConnect} & \textbf{Esparsidad BoostDropout} \\
\thesisimage[width=0.49\textwidth,height=0.78\textheight,keepaspectratio]{imagenes_resultados/autoencoder/DropConnect/first_layer_parameter_sparsity.jpg}{6.0cm} &
\thesisimage[width=0.49\textwidth,height=0.78\textheight,keepaspectratio]{imagenes_resultados/autoencoder/BoostDropout/first_layer_parameter_sparsity.jpg}{6.0cm}
\end{tabular}
\smallsource{imagenes_resultados/autoencoder/DropConnect/first_layer_parameter_sparsity.jpg; imagenes_resultados/autoencoder/BoostDropout/first_layer_parameter_sparsity.jpg}
\end{frame}

% ------------------------------------------------------------
\section{Alcances, conclusiones y cierre}

% ------------------------------------------------------------
\begin{frame}{Conclusiones Generales}
\large
\begin{enumerate}
  \item El sobreajuste aparece como una limitación central en redes neuronales con alta capacidad expresiva.

  \vspace{0.15cm}

  \item Dropout y DropConnect reducen este problema introduciendo dilución aleatoria durante el entrenamiento.

  \vspace{0.15cm}

  \item En esta tesis se propuso \bd{}, una alternativa que reemplaza la supresión binaria por una modulación multiplicativa de las activaciones.

  \vspace{0.15cm}

  \item En el protocolo experimental estudiado, \bdn{} obtuvo menor error y menor pérdida de test que la red base, Dropout y DropConnect.

  \vspace{0.15cm}

  \item La contribución principal es mostrar que la regularización estocástica puede pensarse más allá de ``apagar'' neuronas o conexiones.
\end{enumerate}

\note{
\begin{itemize}
  \item Esta diapositiva debe cerrar el argumento principal de la defensa.
  \item Primero recordar el problema: el sobreajuste.
  \item Luego ubicar Dropout y DropConnect como antecedentes naturales.
  \item Después presentar BoostDropout como extensión conceptual: no anula, sino que modula.
  \item Remarcar que la mejora se observa dentro del protocolo experimental definido.
  \item Evitar presentar el resultado como una superioridad universal.
  \item La frase final es el mensaje más importante: la tesis abre una forma más flexible de pensar la regularización estocástica.
\end{itemize}
}
\end{frame}

% ------------------------------------------------------------
\begin{frame}{Próximos pasos}
\begin{itemize}
  \item Evaluar \bd{} en otros conjuntos de datos.
  \item Probar arquitecturas más profundas y convolucionales.
  \item Explorar con mayor detalle el espacio de $p$, $\lambda$ y normalización.
  \item Analizar calibración, estabilidad y dinámica de entrenamiento.
  \item Estudiar variantes de máscaras multiplicativas más generales.
\end{itemize}

\vspace{0.45cm}

\begin{alertblock}{Mensaje final}
\centering
\large
\textbf{\bd{} no cierra el problema de la regularización: propone una dirección para explorarlo.}
\end{alertblock}

\note{
\begin{itemize}
  \item Esta diapositiva muestra cómo podría continuar la línea abierta por la tesis.
  \item Primero, evaluar si la mejora se mantiene en datasets más variados y exigentes.
  \item Segundo, probar \bd{} en arquitecturas más cercanas al uso actual, como redes convolucionales o más profundas.
  \item Tercero, estudiar de forma más sistemática los hiperparámetros $p$, $\lambda$ y el uso de normalización.
  \item También sería importante analizar calibración, estabilidad y dinámica de entrenamiento para entender mejor el mecanismo.
  \item La frase final funciona como cierre oral: \bd{} no se presenta como solución definitiva, sino como una dirección nueva para explorar. Se puee explorar modulacion multiplicativa de las conexiones por ejemplo.
\end{itemize}
}
\end{frame}

% ------------------------------------------------------------
\begin{frame}[plain]
\centering
\vfill
{\fontsize{48}{58}\bfseries ¡Muchas gracias!}
\vfill
\end{frame}



% ------------------------------------------------------------
%\begin{frame}{Mensaje final}
%\centering
%\vspace{0.4cm}
%{\Large\bfseries Esta tesis propone pensar la regularización no solo como supresión aleatoria, sino como una intervención más general sobre la intensidad de las señales internas de una red neuronal.}
%\vspace{0.8cm}
%
%\thesisimage[width=0.42\linewidth]{14_boostdropout_training.jpg}{3.0cm}
%\vfill
%{\large Muchas gracias.}
%\end{frame}
%
%% ------------------------------------------------------------
%\appendix
%\section{Apéndice}
%
%\begin{frame}{Apéndice: ecuaciones clave de \bd}
%\begin{align*}
% m_j &= \begin{cases}
% 1+\lambda, & \text{con probabilidad } p,\\
% 1, & \text{con probabilidad } 1-p,
% \end{cases} \\
% \widetilde h_j &= m_j h_j, \\
% \mathbb{E}[m_j] &= 1+p\lambda, \\
% \widehat m_j &= \frac{m_j}{1+p\lambda} \quad \Rightarrow \quad \mathbb{E}[\widehat m_j]=1.
%\end{align*}
%\begin{block}{Caso límite}
%Para $\lambda=-1$, la máscara se vuelve binaria; con normalización adecuada se recupera la lógica de \textit{inverted dropout}.
%\end{block}
%\end{frame}
%
%\begin{frame}{Apéndice: tablas completas de test}
%\begin{columns}[T,totalwidth=\textwidth]
%\column{0.49\textwidth}
%\textbf{MMV}
%\begin{table}
%\centering
%\scriptsize
%\begin{tabular}{lccccc}
%\toprule
%Modelo & Media & Mediana & Desv. & Mín. & Máx. \\
%\midrule
%\ofn & .0687 & .0681 & .0023 & .0655 & .0726 \\
%\don & .0609 & .0616 & .0036 & .0547 & .0672 \\
%\dcn & .0614 & .0609 & .0029 & .0555 & .0655 \\
%\bdn & .0480 & .0468 & .0036 & .0442 & .0572 \\
%\bottomrule
%\end{tabular}
%\end{table}
%\column{0.49\textwidth}
%\textbf{Early stopping}
%\begin{table}
%\centering
%\scriptsize
%\begin{tabular}{lccccc}
%\toprule
%Modelo & Media & Mediana & Desv. & Mín. & Máx. \\
%\midrule
%\ofn & .0687 & .0681 & .0023 & .0655 & .0726 \\
%\don & .0609 & .0616 & .0036 & .0547 & .0672 \\
%\dcn & .0614 & .0609 & .0029 & .0555 & .0655 \\
%\bdn & .0491 & .0491 & .0028 & .0451 & .0527 \\
%\bottomrule
%\end{tabular}
%\end{table}
%\end{columns}
%\end{frame}

\end{document}

